\subsection{Early central heights}
\label{sec:early-central-heights}

The first part of the height interval is controlled by exterior
contraction. Besides heights above a coefficient-dependent threshold,
it includes smaller central heights whose positive half-orbits still
have small normalized radius at that threshold. The estimates below
keep the tail error independent of the threshold, so that the common
fixed sections suffice for later applications. We use the scales
\begin{equation}
\label{eq:direct-quadratic-prefix-scales}
 q_0=-a\delta\ell^2,\qquad q_1=b\delta\ell,\qquad
 q_2=-\bar c\delta,\qquad \epsilon_3=\ell/\delta\le1,
\end{equation}
where \(a\ge a_{\min}>0,\ b\ge0,\ \bar c\) lie in a fixed compact
coefficient set and \(\delta>0\) is sufficiently small. Put
\(Q=\delta\ell^2\). In \(x=|\eta|/\ell\) the outward drifts are
\[
 K_\sigma(x)=-q(\sigma\ell x)/(\delta\ell^2)
 =a-\sigma bx+\bar c x^2-\sigma\epsilon_3 x^3R_q(\sigma\ell x).
\]

\begin{lemma}[Early central heights]
\label{lem:early-central-heights}
Fix \(d_f>0\) and put \(z_e=d_fQ\). There is \(r_0>0\), depending only
on \(d_f\) and the compact coefficient set, such that for all
sufficiently small \(\delta\) the central heights in \(\cH_p\) with
\begin{equation}
\label{eq:direct-primary-early}
 h\ge z_e,\qquad\hbox{or}\qquad
 0<h<z_e,\quad |\eta_+(z_e;h)|/\ell\le r_0,
\end{equation}
form an initial interval in every logarithmic height coordinate
\(T=-c_T\log h\), and
\(\Pi_p'(I_-)I_{-,h}/I_{+,h}<1\) on it. In particular
\(\mathcal L<0\) there. The coefficient of \(\sqrt{z_*}\) in the
estimates is \(2C_{\rm cone}\); the constants that depend on
\(d_f\) and \(r_0\) enter only the final threshold for \(\delta\).
\end{lemma}

\begin{proof}
By \cref{lem:central-height-order}, the set
\eqref{eq:direct-primary-early} is an initial interval in \(T\).
There are three ranges: large heights, the intermediate height band,
and smaller heights with small radius at the prescribed cut.
The large-height estimate handles the first. In the other two, we bound the paired
logarithm \(\log(I_{-,h}/I_{+,h})\) of
\cref{lem:paired-variation-identity-new}. The common-cone tail
contributes \(2C_{\rm cone}\sqrt{z_*}\), already below a fixed
fraction of the exterior contraction margin. The remaining errors are
controlled by choosing the radius cut small where needed, and then
reducing \(\delta\).

\emph{Large heights.} Let \(\theta_0\) be the constant of
\cref{lem:relative-height-paired-new}, and fix \(K\ge2a_{\max}/\theta_0\).
For \(h\ge KQ\),
\[
 |q_0|/h\le a/K,\qquad q_1/\sqrt h\le b\sqrt{\delta/K},\qquad
 |q_2|\le C\delta,
\]
so \cref{lem:relative-height-paired-new} gives
\(0<L_p'<1\) for small \(\delta\).

\emph{The band \(d_fQ\le h\le KQ\).} Set \(z=Q\zeta\), \(U=Q\mathcal U\).
Since \(q_0<0\), the constant term of the graph equation is
favorable, and the centered lower barrier \(U_\sigma\ge m(z-h)\) holds
for a small fixed \(m\). On a fixed interval ending at
\(\zeta=K_1>2K\), the graph equation converges uniformly, with its
finite variations, to
\[
 \mathcal U_\zeta=B_0+a/\zeta,\qquad
 \mathcal U(\zeta)=B_0(\zeta-h/Q)+a\log(\zeta Q/h).
\]
Choose \(K_1\) large after \(d_f,K\). Since
\(\log(K_1/d_f)/K_1\to0\), the energy ratio at \(K_1Q\) lies in the
interior of the common cone, uniformly for \(d_f\le h/Q\le K\).
Before this height, the positive flux and the centered lower barrier
bound the three paired integrals by
\(C_{d_f,K_1}(\sqrt Q+q_1/\sqrt Q)=C_{d_f,K_1}(\sqrt Q+b\sqrt\delta)\).
Their square-root singularities at \(\zeta=h/Q\) are integrable.
Beyond this height, \cref{lem:direct-weighted-tail} applies on the
common cone. Combining the two estimates gives
\begin{equation}
\label{eq:direct-primary-early-uniform-section}
 \left|\log(I_{-,h}/I_{+,h})\right|
 \le 2C_{\rm cone}\sqrt{z_*}+C_{d_f,K_1}\sqrt\delta .
\end{equation}

\emph{Small radius at \(z_e\).} Write
\(x_\sigma=|\eta_\sigma(z_e;h)|/\ell\). Take a fixed large \(L_0\), and
then \(r_0\) so small that both \(K_\sigma\) lie between \(a/2\) and
\(2a\) on \(0\le x\le L_0r_0\). Up to a hypothetical first hit
of \(L_0r_0\) by the negative branch, the graph equations give
\(\mathcal G_{-,p}(U_-,z)\le C\mathcal G_{+,p}(U_+,z)\): compare
\(\sup B+2a\delta\ell^2/z\) with \(b_0+a\delta\ell^2/(2z)\), the
supremum being taken on the fixed box. Integration from the common zero
energy gives \(U_-\le CU_+\), and \(L_0^2>C\) excludes that hit.
Thus \(x_+\le r_0\) implies \(x_-\le L_0r_0\).

On the prefix up to \(z_e\), use
\(dz/d|\eta|=|\eta|z/(zB-q)\). The lower bounds
\(zB-q\ge a_{\min}\delta\ell^2/2\) and \(zB-q\ge b_0z\) give
\[
 \int_h^{z_e}\frac{dz}{|\eta_\sigma|}\le C\ell r_0,\qquad
 q_1\int_h^{z_e}\frac{dz}{|\eta_\sigma|z}\le Cbr_0,\qquad
 \int_h^{z_e}\frac{|\eta_\sigma|}{z}\,dz\le C\epsilon_3 r_0^3 .
\]
From \(z_e\) onward there are barriers
\begin{equation}
\label{eq:direct-primary-shifted-cones}
 m(z-z_e)\le U_\sigma(z)
               \le U_\sigma(z_e)+M(z-z_e).
\end{equation}
The initial adverse ratios
\(q_1\sqrt{U_\sigma(z_e)}/z_e\), \(|q_2|U_\sigma(z_e)/z_e\) and
\(U_\sigma(z_e)^{3/2}/z_e\) are \(O(r_0/d_f)\), \(O(r_0^2/d_f)\) and
\(O(\epsilon_3 r_0^3/d_f)\). Choose \(r_0\) small after \(d_f\).
Use the strip between the barriers only up to \(t=\ell^2Z_1\),
with fixed \(Z_1\). Its added linear, quadratic and cubic losses are
\(O(\sqrt\delta)\), \(O(\delta M)\) and \(O(M^{3/2}\ell\sqrt{Z_1})\),
so small \(\delta\) gives positive flux.

These shifted barriers preserve positive flux, but the endpoint
estimate also requires entry into the common cone. For this, apply
the continuation part of
\cref{lem:direct-small-height-continuation} with \(\zeta_0=d_f\) and
\(0\le \mathcal U_0=U_\sigma(z_e)/\ell^2\le Cr_0^2\).
The preceding radius bounds give
\(K_\sigma(\sqrt{2\mathcal U_0})\ge a_{\min}/2\).
Choose the height \(Z_1\) for entry into the cone large after
\(r_0,d_f\). The shifted barriers
bound the paired integrals up to this height by
\(C_{d_f,r_0,Z_1}\{\sqrt\delta+\ell(1+|\log\delta|)\}\):
the linear integral is \(O(\sqrt\delta)\), and splitting the square
root in the upper barrier bounds the cubic integral. This tends to zero
uniformly for \(0<\epsilon_3\le1\), since \(\ell=\epsilon_3\delta\).
The remaining common-cone tail contributes
\(2C_{\rm cone}\sqrt{z_*}+O(\delta)\). Together with the prefix
bounds, this proves contraction for small \(r_0\) and then small
\(\delta\).
\end{proof}
